\chapter{Lactate (Lactic Acid)}

\section*{What Is This Test?}
The lactate blood test measures the concentration of lactic acid (lactate) in the blood. Lactate is the end product of anaerobic glycolysis, produced when pyruvate (the product of glycolysis) cannot enter the mitochondria for oxidative phosphorylation due to insufficient oxygen delivery or impaired mitochondrial function. The interconversion of pyruvate and lactate is catalyzed by lactate dehydrogenase (LDH): Pyruvate + NADH + H\textsuperscript{+} $\leftrightarrow$ Lactate + NAD\textsuperscript{+}. Under normal aerobic conditions, the lactate/pyruvate ratio is approximately 10:1, with serum lactate maintained at <2 mmol/L. Lactate is produced by all tissues but especially by skeletal muscle, red blood cells, brain, and skin. It is cleared primarily by the liver (60--70\%, via the Cori cycle---lactate converted back to glucose via gluconeogenesis) and kidneys (30\%, excreted or metabolized). Lactic acidosis is defined as serum lactate >4--5 mmol/L with metabolic acidosis (pH <7.35) \cite{kraut2014lactic}. Lactic acidosis is classified as Type A (tissue hypoxia) or Type B (impaired cellular metabolism without clinical evidence of global hypoxia). Serum lactate is a critical biomarker in sepsis, shock, trauma, and critical illness. Lactate >4 mmol/L in sepsis is associated with 28--38\% mortality, and lactate clearance (a decrease in lactate by >10--20\% over 2--6 hours with resuscitation) is a marker of adequate treatment response. The Surviving Sepsis Campaign recommends measuring lactate within 1 hour of sepsis recognition, with repeat measurement within 2--4 hours if initial lactate is >2 mmol/L \cite{evans2021surviving,tietz2023textbook}.

\section*{Alternative Names}
\begin{itemize}
  \item Lactic Acid
  \item Serum Lactate
  \item Blood Lactate
  \item Plasma Lactate
  \item L-Lactate
  \item D-Lactate
  \item Lactate, Whole Blood
\end{itemize}
\section*{Principle}
Lactate is measured by an enzymatic method using lactate oxidase: L-lactate + O\textsubscript{2} $\rightarrow$ pyruvate + H\textsubscript{2}O\textsubscript{2} (lactate oxidase). The H\textsubscript{2}O\textsubscript{2} produced is detected by a peroxidase-coupled reaction: H\textsubscript{2}O\textsubscript{2} + chromogen $\rightarrow$ colored product + H\textsubscript{2}O. The color intensity is proportional to lactate concentration (absorbance at 540--600 nm). Alternatively, some analyzers use an amperometric electrode that directly measures the H\textsubscript{2}O\textsubscript{2} current. The lactate oxidase method is specific for L-lactate, the physiological isomer. D-lactate (produced by gut bacteria in short bowel syndrome) requires a separate D-lactate assay using D-lactate dehydrogenase. Most point-of-care lactate analyzers (i-STAT, blood gas analyzers) use the amperometric lactate oxidase method. On automated chemistry analyzers, the spectrophotometric method is standard. Critical pre-analytical requirement: Lactate must be measured rapidly after collection because red blood cells continue to produce lactate through glycolysis in vitro. Either the sample must be analyzed within 15 minutes, or it must be collected in a sodium fluoride/potassium oxalate tube (gray-top) which inhibits glycolysis, or it must be placed on ice immediately.

\section*{History}
Lactic acid was first isolated from sour milk by Carl Wilhelm Scheele in 1780. The role of lactate in muscle metabolism was described by Otto Meyerhof and Archibald Hill in the 1920s, for which they received the Nobel Prize in 1922. The Cori cycle was described by Carl and Gerty Cori in the 1920s-1930s. The clinical significance of lactate in shock was established in the 1960s by Weil, Broder, and others. The use of serial lactate measurements to guide resuscitation in sepsis was established by Rivers and colleagues in the landmark early goal-directed therapy (EGDT) trial (2001) \cite{rivers2001early} and subsequently refined. Lactate-guided resuscitation is now a cornerstone of sepsis management per the Surviving Sepsis Campaign guidelines.

\section*{Why This Test Is Ordered}
\begin{itemize}
  \item Diagnosis and severity assessment of sepsis and septic shock: lactate >2 mmol/L is one criterion for sepsis; lactate >4 mmol/L defines septic shock and predicts mortality
  \item Monitoring of resuscitation adequacy in shock states: serial lactate measurements every 2--4 hours to assess lactate clearance. Failure to clear lactate to normal (<2 mmol/L) within 6--24 hours is associated with increased mortality
  \item Evaluation of unexplained high-anion-gap metabolic acidosis: lactate is the most common cause of elevated anion gap metabolic acidosis in hospitalized patients
  \item Diagnosis of mesenteric ischemia (acute intestinal ischemia): markedly elevated lactate (>5--10 mmol/L) with severe abdominal pain out of proportion to examination; requires emergent CT angiography or surgical exploration
  \item Assessment of tissue perfusion in cardiac arrest, major trauma, burns, and post-surgical patients
  \item Detection of metformin-associated lactic acidosis (MALA): rare but life-threatening complication of metformin therapy, especially with renal impairment, hepatic disease, or acute illness; lactate >10 mmol/L with severe acidosis
  \item Diagnosis of D-lactic acidosis in patients with short bowel syndrome presenting with episodic metabolic acidosis and neurological symptoms (confusion, ataxia, slurred speech)
  \item Evaluation of suspected mitochondrial disorders, inborn errors of metabolism, or unexplained lactic acidosis
\end{itemize}

\section*{Primary vs Secondary Test}
Lactate is a primary test in critical care and emergency medicine. It is not part of routine panels (BMP/CMP) and must be specifically ordered. It is a core biomarker in the Surviving Sepsis Campaign bundle and is essential in any patient with suspected shock or tissue hypoperfusion.

\section*{How This Test Is Performed}
Lactate is measured from venous blood (most common), arterial blood (for concurrent blood gas analysis), or capillary blood (point-of-care). The sample MUST be handled correctly: (1) Use a gray-top tube (sodium fluoride/potassium oxalate) which inhibits glycolysis and prevents in vitro lactate production by red blood cells. A green-top (lithium heparin) tube can be used only if analyzed within 15 minutes. (2) If a gray-top tube is not available, the sample must be placed on ice immediately and analyzed within 30 minutes. (3) The tourniquet should be released promptly---prolonged venous stasis causes local tissue hypoxia and lactate accumulation, falsely elevating lactate from the drawn limb by up to 1--2 mmol/L. (4) The sample should not be drawn from an arm receiving IV fluids containing lactate (lactated Ringer's) or glucose. Results from point-of-care devices (i-STAT, blood gas analyzers) are available within 2--5 minutes; central laboratory results are available within 30--60 minutes.

\section*{What Preparation Is Needed}
No fasting is required, but the patient should avoid vigorous exercise for at least 1 hour before testing (exercise lactate can exceed 10 mmol/L and take 30--60 minutes to clear). Avoid fist clenching or prolonged tourniquet application during venipuncture (local stasis falsely elevates lactate). Inform the provider of: (1) Metformin use (risk of MALA, especially with renal impairment or acute illness). (2) Antiretroviral therapy (NRTIs, especially stavudine and didanosine, can cause lactic acidosis from mitochondrial toxicity). (3) Propofol infusion syndrome (prolonged high-dose propofol can cause lactic acidosis). (4) Linezolid (prolonged use can cause lactic acidosis from mitochondrial toxicity). (5) Epinephrine and other catecholamine infusions (beta-2 stimulation increases glycolysis and lactate production). (6) Alcohol intoxication (increases lactate from ethanol metabolism generating NADH, shifting pyruvate to lactate).

\section*{Contraindications and Precautions}
\begin{itemize}
  \item No absolute contraindications. Critical pre-analytical considerations:
  \item (1) In vitro glycolysis: red blood cells continue to produce lactate after blood is drawn. Without a glycolysis inhibitor (fluoride tube) or rapid processing on ice, lactate can increase by 0.5--1.0 mmol/L per 30 minutes at room temperature.
  \item (2) Prolonged tourniquet time with fist clenching: local stasis hypoxia causes lactate accumulation; can elevate measured lactate by up to 1--2 mmol/L.
  \item (3) Venous vs. arterial: arterial lactate is slightly lower than venous; capillary (fingerstick) lactate correlates well with venous.
  \item (4) Delayed processing: if the sample sits at room temperature without fluoride inhibition, lactate will be falsely elevated.
  \item (5) Epinephrine, albuterol, and other beta-agonists cause a transient rise in lactate by stimulating glycolysis (Type B lactic acidosis---often 3--6 mmol/L, has no adverse prognostic significance, and resolves when the drug is discontinued).
  \item (6) Seizures: lactate rises dramatically during generalized tonic-clonic seizures (can exceed 10--15 mmol/L) but clears rapidly (within 60--90 minutes) post-ictally.
  \item (7) D-lactic acidosis is not detected by standard L-lactate assays; D-lactate must be specifically ordered.
  \item (8) Hemolysis does not significantly affect lactate measurement with the lactate oxidase method.
\end{itemize}
\section*{Risks}
Minimal risk from venipuncture. Mild pain or bruising resolves in 1--2 days. Arterial puncture for ABG has slightly higher risk (hematoma, arterial dissection, rare distal ischemia).

\section*{What Happens After the Test}
Results available within 5--60 minutes depending on testing location. Critical value: lactate >4 mmol/L triggers the sepsis bundle (blood cultures, antibiotics within 1 hour, IV fluids 30 mL/kg, vasopressors if hypotensive, remeasure lactate within 2--4 hours) \cite{tierney2023current}. Lactate clearance <10\% at 2--4 hours after initial resuscitation is associated with increased mortality and should prompt reassessment of resuscitation adequacy. Lactate elevation in sepsis is multifactorial: tissue hypoxia from hypoperfusion (Type A), but also from catecholamine-driven accelerated glycolysis, impaired hepatic clearance, mitochondrial dysfunction (Type B). Thus, lactate is a marker of illness severity and metabolic stress, not purely a marker of hypoxia \cite{suetrong2016lactic}. Trending lactate over time is more informative than a single measurement.

\section*{Understanding Results}

\subsection*{Below Normal}
\begin{itemize}
  \item Low lactate (<0.5 mmol/L) is not clinically significant. Normally insignificant.
  \item May be seen in severe liver failure (reduced lactate production due to impaired Cori cycle) but this is not clinically useful.
  \item In rare inborn errors of gluconeogenesis, lactate may be paradoxically low due to substrate deficiency.
\end{itemize}

\subsection*{Above Normal}
\begin{itemize}
  \item \textbf{Lactate >2 mmol/L:} Above normal; requires clinical attention.
  \item \textbf{Lactate >4 mmol/L:} Defines lactic acidosis. Associated with significantly increased mortality in sepsis (28--38\%). Triggers sepsis protocol.
  \item \textbf{Type A lactic acidosis (tissue hypoxia):} (a) Shock: septic, cardiogenic, hypovolemic, obstructive, distributive. Lactate is the central biomarker in sepsis---measures tissue hypoperfusion and metabolic stress. (b) Cardiac arrest: lactate rises during arrest due to global ischemia; failure to clear lactate post-ROSC (return of spontaneous circulation) indicates ongoing hypoperfusion or severe tissue injury. (c) Severe hypoxemia: respiratory failure, ARDS, carbon monoxide poisoning, cyanide poisoning. (d) Severe anemia: oxygen-carrying capacity severely reduced. (e) Mesenteric ischemia: acute intestinal ischemia---lactate often >5 mmol/L, may be the earliest lab abnormality. (f) Regional hypoperfusion: compartment syndrome, limb ischemia, massive trauma. (g) Strenuous exercise: transient lactate elevation (can exceed 10 mmol/L) that clears within 30--60 minutes.
  \item \textbf{Type B lactic acidosis (no clinical evidence of tissue hypoxia):} (a) Sepsis (also has Type B component---mitochondrial dysfunction and accelerated glycolysis from inflammatory mediators). (b) Metformin-associated lactic acidosis (MALA): rare (3--9 cases/100,000 patient-years) but life-threatening; occurs with metformin accumulation in renal impairment (eGFR <30), hepatic disease, or acute illness. Lactate typically >10 mmol/L, severe metabolic acidosis, pH <7.0. Treatment: supportive care, hemodialysis removes metformin and corrects acidosis. (c) Malignancy: hematologic malignancies (leukemia, lymphoma) and solid tumors can produce lactate through the Warburg effect (aerobic glycolysis in tumor cells). (d) Liver failure: impaired hepatic clearance of lactate---liver clears 60--70\% of daily lactate; in fulminant hepatic failure and decompensated cirrhosis, lactate accumulates. In acute liver failure, arterial lactate >3.5 mmol/L at 4 hours or >3.0 at 12 hours after presentation is used for transplant listing criteria (King's College criteria). (e) Thiamine deficiency: beriberi causing impaired pyruvate dehydrogenase activity and lactate accumulation (shoshin beriberi presents with fulminant high-output heart failure and severe lactic acidosis). (f) Medications and toxins: metformin, NRTIs (stavudine, didanosine---mitochondrial toxicity), linezolid (mitochondrial toxicity with prolonged use), propofol (propofol infusion syndrome), epinephrine and other beta-agonists (stimulate glycolysis), cyanide, carbon monoxide, methanol, ethylene glycol (metabolites impair oxidative phosphorylation). (g) D-lactic acidosis: short bowel syndrome with carbohydrate malabsorption; unabsorbed carbohydrates are fermented by colonic bacteria to D-lactate, which humans metabolize slowly. Presents with episodic high-anion-gap metabolic acidosis and neurological symptoms (confusion, ataxia, dysarthria) after high-carbohydrate meals. Standard L-lactate assays do NOT detect D-lactate; must be specifically ordered \cite{uribarri1998dlactic}. (h) Seizures: during and immediately post-ictal, lactate can exceed 10--15 mmol/L due to massive muscle activity and local tissue hypoxia. Clears within 60--90 minutes. (i) Diabetic ketoacidosis: mild-to-moderate lactate elevation (3--5 mmol/L) from volume depletion, tissue hypoperfusion, and accelerated glycolysis. (j) Alcoholic ketoacidosis: similar mechanism.
  \item \textbf{Lactate clearance:} (Initial lactate -- repeat lactate) / initial lactate $\times$ 100\%. Clearance >10--20\% at 2--4 hours indicates effective resuscitation. Failure to clear lactate (<10\%) portends poor outcome.
\end{itemize}

\section*{Normal Results}
\begin{itemize}
  \item \textbf{Venous lactate (adults):} <2 mmol/L (0.5--2.0 mmol/L or 4.5--18 mg/dL) \cite{wallach2022interpretation}
  \item \textbf{Arterial lactate:} <1.5 mmol/L (slightly lower than venous)
  \item \textbf{Critical value (sepsis):} $\geq$4 mmol/L (requires immediate intervention per Surviving Sepsis Campaign)
  \item \textbf{Capillary/fingerstick lactate (point-of-care):} <2 mmol/L
  \item \textbf{Conditions:} No tourniquet stasis, no fist clenching, no recent exercise
  \item \textbf{Conversion:} 1 mmol/L = 9 mg/dL
\end{itemize}

\section*{Follow-up Tests}
\begin{itemize}
  \item \textbf{Repeat lactate (serial measurements):} Every 2--4 hours in sepsis/shock to assess resuscitation adequacy and lactate clearance.
  \item \textbf{Arterial blood gas (ABG):} pH, pCO\textsubscript{2}, bicarbonate, and base deficit---defines acidosis severity and changes.
  \item \textbf{Anion gap:} Elevated AG suggests lactic acidosis or other unmeasured acid anions.
  \item \textbf{Comprehensive metabolic panel (CMP):} Renal function (BUN, creatinine), liver function, glucose.
  \item \textbf{Blood cultures, lactate, and early antibiotics:} Core sepsis bundle (Surviving Sepsis Campaign).
  \item \textbf{D-lactate (if short bowel syndrome suspected):} Specific D-lactate assay.
  \item \textbf{Creatine kinase (CK):} If seizure or rhabdomyolysis is cause of lactate elevation.
  \item \textbf{Acetaminophen level, salicylate level, toxic alcohol screen:} For high-AG metabolic acidosis with elevated lactate.
  \item \textbf{Metformin level:} If metformin-associated lactic acidosis suspected.
\end{itemize}

\section*{References}
\bibliographystyle{unsrtnat}
\bibliography{references}

\clearpage
\section*{Regulatory Status and Authorization Date}
This chapter describes a test category, not a specific commercial product. FDA, CDSCO, EU IVDR, and MHRA approval or registration dates therefore cannot be assigned without the assay manufacturer, product name, instrument, and intended use. For laboratory-developed tests, an individual agency approval date may not exist. See the Preface for the interpretation of regulatory status.

\clearpage
